\subsection{Nonparametric Curve Estimation}
% Density and regression curves, smoothing parameters, histogram and kernel estimators.

Nonparametric curve estimation studies unknown functions.\\
Unlike empirical distribution functions, these targets cannot be estimated well at fine resolution without smoothness. The practical question is therefore how much to smooth the data.

\begin{definition} \label{def:curve-estimator-ise-mise} \hlt{Curve Estimator, ISE, and MISE}\\
Let $g$ be an unknown function and let $\widehat g_n$ be a sample-dependent estimator of $g$.\\
At fixed $x$, $\widehat g_n(x)$ is random, while $x$ is deterministic.
\begin{enumerate}[label=(\roman*)]
\setlength{\itemsep}{0pt}
\item The \hlt{integrated squared error} is $L(g,\widehat g_n)\defeq\int (g(u)-\widehat g_n(u))^2\,du$.
\item The \hlt{risk}, or \hlt{mean integrated squared error}, is $R(g,\widehat g_n)\defeq\mathbb{E}[L(g,\widehat g_n)]$.
\end{enumerate}
\end{definition}

\begin{theorem} \label{thm:curve-risk-bias-variance-decomposition} \hlt{Integrated Bias--Variance Decomposition}\\
Assume $R(g,\widehat g_n)<\infty$ and expectation may be interchanged with integration.\\
Define $b(x)\defeq\mathbb{E}\widehat g_n(x)-g(x)$ and $v(x)\defeq\operatorname{Var}(\widehat g_n(x))$. Then $R(g,\widehat g_n)=\int b(x)^2\,dx+\int v(x)\,dx$.
\begin{proof} For fixed $x$, put $\mu(x)\defeq\mathbb{E}\widehat g_n(x)$.
Since $g(x)$ is deterministic, $\mathbb{E}[(\widehat g_n(x)-g(x))^2]=\mathbb{E}[(\widehat g_n(x)-\mu(x))^2]+(\mu(x)-g(x))^2=v(x)+b(x)^2$.\\
Hence $R(g,\widehat g_n)=\mathbb{E}\int(\widehat g_n(x)-g(x))^2\,dx=\int\mathbb{E}[(\widehat g_n(x)-g(x))^2]\,dx=\int v(x)\,dx+\int b(x)^2\,dx$.\\
\hlt{Structural Commentary}: Integrated risk is the pointwise mean-squared error identity integrated over $x$.\\
\hlt{Key Insight}: Variance measures estimator fluctuation; squared bias measures systematic displacement.
\end{proof}
\end{theorem}

\begin{figure}[H]
  \centering
  \begin{tikzpicture}[
  x=1cm,
  y=1cm,
  >=Stealth,
  line cap=round,
  line join=round,
  every node/.style={font=\small},
  axis/.style={->, black!70, line width=0.8pt},
  risk/.style={black!78, line width=0.95pt},
  bias/.style={red!70!black, line width=0.9pt},
  variance/.style={green!50!black, densely dashed, line width=0.9pt},
  guide/.style={black!55, densely dashed, line width=0.65pt}
]
  \draw[axis] (0,0) -- (6.05,0) node[right] {smoothing};
  \draw[axis] (0,0) -- (0,3.25) node[above] {error};

  \draw[bias]
    plot[smooth] coordinates {(0.45,0.15) (1.2,0.22) (2.0,0.42) (2.8,0.75) (3.7,1.22) (4.6,1.82) (5.35,2.55)}
    node[right] {$\mathrm{bias}^2$};
  \draw[variance]
    plot[smooth] coordinates {(0.45,2.75) (1.2,2.05) (2.0,1.50) (2.8,1.10) (3.7,0.85) (4.6,0.70) (5.35,0.62)}
    node[right] {variance};
  \draw[risk]
    plot[smooth] coordinates {(0.45,2.90) (1.2,2.25) (2.0,1.86) (2.8,1.72) (3.7,1.88) (4.6,2.35) (5.35,3.02)}
    node[right] {risk};

  \draw[guide] (2.8,0) -- (2.8,1.72);
  \filldraw[white, draw=black!75, line width=0.65pt] (2.8,1.72) circle (1.45pt);
  \node[below] at (2.8,0) {balance};
\end{tikzpicture}

  \caption{Smoothing trades variance for bias; risk is minimized near balance.}
  \label{fig:bias-variance-smoothing-tradeoff}
\end{figure}

\begin{principle} \label{princ:curve-estimation-smoothing-tradeoff} \hlt{Smoothing Trades Variance for Bias}\\
Increasing a smoothing parameter averages over more observations, decreasing $\int v(x)\,dx$ but increasing $\int b(x)^2\,dx$ by erasing local structure.
Smoothing is chosen by balancing these two risk terms.
\end{principle}

\subsubsection{Histograms}

\begin{definition} \label{def:regular-histogram-density-estimator} \hlt{Regular Histogram Estimator}\\
Let $X_1,\ldots,X_n$ be IID from density $f$ on $[0,1]$.\\
For $m\geq1$, set $h=1/m$, $B_j=[(j-1)h,jh)$ for $j<m$, and $B_m=[(m-1)h,1]$.\\
Let $\nu_j\defeq\sum_{i=1}^{n}\mathbf{1}_{\{X_i\in B_j\}}$, $\widehat p_j\defeq\nu_j/n$, and $p_j\defeq\int_{B_j}f(u)\,du$.\\
The \hlt{histogram density estimator} is $\widehat f_n(x)\defeq\sum_{j=1}^{m}(\widehat p_j/h)\mathbf{1}_{\{x\in B_j\}}$.
\end{definition}

\begin{proposition} \label{prop:histogram-pointwise-mean-variance} \hlt{Histogram Mean and Variance}\\
Fix $m$ and $x\in B_j$. Then $\mathbb{E}\widehat f_n(x)=p_j/h$ and $\operatorname{Var}(\widehat f_n(x))=p_j(1-p_j)/(nh^2)$.
\begin{proof} Since $\nu_j=\sum_{i=1}^{n}\mathbf{1}_{\{X_i\in B_j\}}$, $\nu_j\sim\operatorname{Binomial}(n,p_j)$. Thus $\mathbb{E}\widehat p_j=p_j$ and $\operatorname{Var}(\widehat p_j)=p_j(1-p_j)/n$.\\
For $x\in B_j$, $\widehat f_n(x)=\widehat p_j/h$, so scaling gives the displayed mean and variance.\\
\hlt{Structural Commentary}: A histogram height is a rescaled bin proportion.\\
\hlt{Key Insight}: Binomial count variance is amplified by the factor $h^{-2}$.
\end{proof}
\end{proposition}

\begin{theorem} \label{thm:histogram-risk-and-optimal-binwidth} \hlt{Histogram Risk and Optimal Binwidth}\\
Let $A\defeq\int_0^1(f'(u))^2\,du\in(0,\infty)$. For small $h$ and large $n$, $R(f,\widehat f_n)\approx Ah^2/12+1/(nh)$.\\
The minimizing binwidth is $h^*=(6/(nA))^{1/3}$, and the corresponding risk is of order $n^{-2/3}$.
\begin{proof} For $x\in B_j$, $p_j/h$ is the bin average of $f$.\\
By Taylor expansion on bin midpoint $c_j$, $p_j/h-f(x)\approx f'(c_j)(c_j-x)$, hence $\int_{B_j}(p_j/h-f(x))^2\,dx\approx f'(c_j)^2h^3/12$.\\
Summing over $j$ gives integrated squared bias $\approx Ah^2/12$.\\
Also $p_j\approx f(c_j)h$ and $1-p_j\approx1$, so $\int_{B_j}\operatorname{Var}(\widehat f_n(x))\,dx\approx h\,p_j/(nh^2)=p_j/(nh)$, and summing gives $\approx1/(nh)$.
Thus the risk proxy is $Q(h)=Ah^2/12+1/(nh)$.\\
Solving $Q'(h)=Ah/6-1/(nh^2)=0$ gives $h^*=(6/(nA))^{1/3}$, and $Q(h^*)=O(n^{-2/3})$.\\
\hlt{Structural Commentary}: Histogram risk is approximated by squared bin-averaging bias plus binomial variance.\\
\hlt{Key Insight}: The optimal $h$ balances $h^2$ bias against $(nh)^{-1}$ variance.
\end{proof}
\end{theorem}

\begin{method} \label{meth:histogram-cross-validation} \hlt{Choosing Histogram Binwidth by Cross-Validation}
\begin{enumerate}[label=(\roman*)]
\setlength{\itemsep}{0pt}
\item Since $\int f^2$ is independent of $h$, minimizing $\mathbb{E}\int(\widehat f_n-f)^2$ is equivalent to minimizing $\mathbb{E}J(h)$,\\
where $J(h)\defeq\int\widehat f_n(x)^2\,dx-2\int\widehat f_n(x)f(x)\,dx$.
\item Let $\widehat f_{-i}$ be the histogram estimator after deleting $X_i$,\\
and estimate $J(h)$ by $\widehat J(h)\defeq\int\widehat f_n(x)^2\,dx-(2/n)\sum_{i=1}^{n}\widehat f_{-i}(X_i)$.
\item For a regular histogram, $\int\widehat f_n^2=h^{-1}\sum_j\widehat p_j^2$ and $\sum_i\widehat f_{-i}(X_i)=((n-1)h)^{-1}\sum_{j=1}^{m}\nu_j(\nu_j-1)$.
\item Thus $\widehat J(h)=2/((n-1)h)-((n+1)/((n-1)h))\sum_{j=1}^{m}\widehat p_j^2$, and choose $h$ minimizing this score.
\end{enumerate}
\end{method}

\begin{method} \label{meth:histogram-confidence-envelope} \hlt{Histogram Confidence Envelope}
\begin{enumerate}[label=(\roman*)]
\setlength{\itemsep}{0pt}
\item For $x\in B_j$, define the histogramized target $\widetilde f(x)\defeq p_j/h$.
\item Set $c\defeq z_{\alpha/(2m)}\sqrt{m/n}/2$.
\item Define $\ell(x)\defeq(\max\{\sqrt{\widehat f_n(x)}-c,0\})^2$ and $u(x)\defeq(\sqrt{\widehat f_n(x)}+c)^2$.
\item If $m=m(n)$ and $m\log n/n\to0$, then $\mathbb{P}(\ell(x)\leq\widetilde f(x)\leq u(x)\ \forall x)\approx1-\alpha$.
\end{enumerate}
\end{method}

\subsubsection{Kernel Density Estimation}

\begin{definition} \label{def:kernel-density-estimator} \hlt{Kernel Density Estimator}\\
\begin{enumerate}[label=(\roman*)]
\setlength{\itemsep}{0pt}
\item A \hlt{kernel} is a smooth $K$ with $K\geq0$, $\int K=1$, $\int xK(x)\,dx=0$, and $\sigma_K^2\defeq\int x^2K(x)\,dx>0$.
\item Given bandwidth $h>0$, the \hlt{kernel density estimator} is $\widehat f_n(x)\defeq n^{-1}\sum_{i=1}^{n}h^{-1}K((x-X_i)/h)$.
\end{enumerate}
Common kernels include the Gaussian $K(x)=(2\pi)^{-1/2}e^{-x^2/2}$ and Epanechnikov $K(x)=\frac34(1-x^2/5)\mathbf{1}_{\{|x|<\sqrt5\}}/\sqrt5$.
\end{definition}

\begin{figure}[H]
  \centering
  \begin{tikzpicture}[
  x=0.55cm,
  y=18cm,
  >=Stealth,
  line cap=round,
  line join=round,
  every node/.style={font=\small},
  axis/.style={->, black!70, line width=0.8pt},
  kernel/.style={red!65!black, line width=0.6pt},
  estimate/.style={black!78, line width=0.95pt},
  sample/.style={green!55!black, line width=0.85pt}
]
  \draw[axis] (-5.2,0) -- (7.1,0) node[right] {$x$};
  \draw[axis] (-5,0) -- (-5,0.070) node[above] {density};

  \foreach \x in {-2.8,-0.8,1.6,3.1,4.2} {
    \draw[sample] (\x,0) -- (\x,0.018);
    \fill[green!55!black] (\x,0.018) circle (1.35pt);
  }

  \draw[kernel]
    plot[smooth] coordinates {(-4.7,0) (-4.1,0.002) (-3.4,0.018) (-2.8,0.038) (-2.2,0.018) (-1.6,0.002) (-1.0,0)};
  \draw[kernel]
    plot[smooth] coordinates {(-2.7,0) (-2.1,0.002) (-1.4,0.018) (-0.8,0.038) (-0.2,0.018) (0.4,0.002) (1.0,0)};
  \draw[kernel]
    plot[smooth] coordinates {(-0.3,0) (0.3,0.002) (1.0,0.018) (1.6,0.038) (2.2,0.018) (2.8,0.002) (3.4,0)};
  \draw[kernel]
    plot[smooth] coordinates {(1.2,0) (1.8,0.002) (2.5,0.018) (3.1,0.038) (3.7,0.018) (4.3,0.002) (4.9,0)};
  \draw[kernel]
    plot[smooth] coordinates {(2.3,0) (2.9,0.002) (3.6,0.018) (4.2,0.038) (4.8,0.018) (5.4,0.002) (6.0,0)};

  \draw[estimate]
    plot[smooth] coordinates {(-4.7,0.0005) (-3.7,0.006) (-2.8,0.034) (-2.0,0.020) (-0.8,0.038) (0.0,0.014)
    (0.9,0.015) (1.6,0.043) (2.3,0.023) (3.1,0.050) (3.7,0.044) (4.2,0.048) (5.0,0.017) (6.0,0.002)};

  \node[below] at (-2.8,0) {\scriptsize $X_1$};
  \node[below] at (-0.8,0) {\scriptsize $X_2$};
  \node[below] at (1.6,0) {\scriptsize $X_3$};
  \node[below] at (3.1,0) {\scriptsize $X_4$};
  \node[below] at (4.2,0) {\scriptsize $X_5$};
\end{tikzpicture}

  \caption{Kernel density estimate: $\widehat f_n(x)=n^{-1}\sum_i h^{-1}K((x-X_i)/h)$.}
  \label{fig:kernel-density-lumps}
\end{figure}

\begin{theorem} \label{thm:kernel-density-risk-rate} \hlt{Kernel Density Risk Rate}\\
Assume $f$ twice differentiable, $c_1\defeq\int u^2K(u)\,du\in(0,\infty)$, $c_2\defeq\int K(u)^2\,du<\infty$, $c_3\defeq\int(f''(x))^2\,dx\in(0,\infty)$.\\
For small $h$ and large $n$, $R(f,\widehat f_n)\approx c_1^2c_3h^4/4+c_2/(nh)$.\\
The minimising bandwidth is $h^*=c_1^{-2/5}c_2^{1/5}c_3^{-1/5}n^{-1/5}$, and the corresponding risk is of order $n^{-4/5}$.
\begin{proof} Let $K_h(x,X)=h^{-1}K((x-X)/h)$, so $\mathbb{E}\widehat f_n(x)=\int K(u)f(x-hu)\,du$, $\operatorname{Var}(\widehat f_n(x))=n^{-1}\operatorname{Var}(K_h(x,X))$.\\
Taylor expansion gives $f(x-hu)=f(x)-hu f'(x)+h^2u^2f''(x)/2+o(h^2)$. Using $\int K=1$ and $\int uK(u)\,du=0$, the pointwise bias is $c_1h^2f''(x)/2+o(h^2)$, so integrated squared bias is $c_1^2c_3h^4/4+o(h^4)$.\\
Also $\mathbb{E}K_h(x,X)^2=h^{-1}\int K(u)^2f(x-hu)\,du\approx h^{-1}c_2f(x)$, hence integrated variance is $\approx c_2/(nh)$.\\
Thus the risk proxy is $Q(h)=c_1^2c_3h^4/4+c_2/(nh)$.\\
Solving $Q'(h)=c_1^2c_3h^3-c_2/(nh^2)=0$ gives $h^*$ and $Q(h^*)=O(n^{-4/5})$.\\
\hlt{Structural Commentary}: Kernel risk is fourth-order squared bias plus first-order inverse effective sample size.\\
\hlt{Key Insight}: The optimal bandwidth balances $h^4$ bias against $(nh)^{-1}$ variance.
\end{proof}
\end{theorem}

\begin{proposition} \label{prop:boxcar-kernel-expectation-variance-consistency} \hlt{Boxcar Kernel Exercise}\\
Let $K(t)=\mathbf{1}_{\{-1/2<t<1/2\}}$, $\widehat f_n(x)\defeq n^{-1}\sum_{i=1}^{n}h^{-1}K((x-X_i)/h)$, and $p_h(x)\defeq\int_{x-h/2}^{x+h/2}f(y)\,dy$.\\
Then $\mathbb{E}\widehat f_n(x)=p_h(x)/h$ and $\operatorname{Var}(\widehat f_n(x))=(p_h(x)-p_h(x)^2)/(nh^2)$.\\
If $f$ is continuous at $x$, $h\to0$, and $nh\to+\infty$, then $\widehat f_n(x)\xrightarrow{\mathbb{P}}f(x)$.
\begin{proof}
Let $N_h(x)\defeq\sum_{i=1}^{n}\mathbf{1}_{\{|X_i-x|<h/2\}}$.
Then $N_h(x)\sim\operatorname{Binomial}(n,p_h(x))$ and $\widehat f_n(x)=N_h(x)/(nh)$.
Thus $\mathbb{E}\widehat f_n(x)=p_h(x)/h$ and $\operatorname{Var}(\widehat f_n(x))=p_h(x)(1-p_h(x))/(nh^2)$.
Continuity gives $p_h(x)/h\to f(x)$, while $p_h(x)=O(h)$ gives $\operatorname{Var}(\widehat f_n(x))=O((nh)^{-1})\to0$.
Chebyshev's inequality yields $\widehat f_n(x)\xrightarrow{\mathbb{P}}f(x)$.\\
\hlt{Structural Commentary}: The boxcar estimator is a rescaled local binomial count.\\
\hlt{Key Insight}: Consistency needs the window to shrink but the expected local count $nh$ to diverge.
\end{proof}
\end{proposition}

\begin{method} \label{meth:kde-cross-validation} \hlt{Kernel Density Cross-Validation}
\begin{enumerate}[label=(\roman*)]
\setlength{\itemsep}{0pt}
\item For bandwidth selection, use $\widehat J(h)\defeq\int\widehat f_n(x)^2\,dx-(2/n)\sum_{i=1}^{n}\widehat f_{-i}(X_i)$, where $\widehat f_{-i}$ omits $X_i$.
\item Define $K^{(2)}(z)\defeq\int K(z-y)K(y)\,dy$ and $K^*(z)\defeq K^{(2)}(z)-2K(z)$.
\item Use the computational shortcut $\widehat J(h)\approx (hn^2)^{-1}\sum_i\sum_j K^*((X_i-X_j)/h)+2K(0)/(nh)$.
\item If $K$ is standard Gaussian, then $K^{(2)}$ is the $N(0,2)$ density.
\end{enumerate}
\end{method}

\begin{theorem} \label{thm:stone-cross-validation-oracle-risk} \hlt{Stone's Cross-Validation Oracle Property}\\
$\mathcal{H}_n$ candidate bandwidth set, $\widehat h_n\in\operatorname*{arg\,min}_{h\in\mathcal{H}_n}\widehat J_n(h)$, $L_n(h)\defeq\int(f(x)-\widehat f_h(x))^2\,dx$, $L_n^*\defeq\inf_{h\in\mathcal{H}_n}L_n(h)>0$.\\
Assume for random $C_n$ not depending on $h$, $\sup_{h\in\mathcal{H}_n}|\widehat J_n(h)-C_n-L_n(h)|=o_{\mathbb{P}}(L_n^*)$. Then $L_n(\widehat h_n)/L_n^*\xrightarrow{\mathbb{P}}1$.
\begin{proof} Let $\Delta_n\defeq\sup_{h\in\mathcal{H}_n}|\widehat J_n(h)-C_n-L_n(h)|$. For any $\eta>0$, choose $h_\eta\in\mathcal{H}_n$ with $L_n(h_\eta)\leq(1+\eta)L_n^*$.\\
Since $\widehat h_n$ minimizes $\widehat J_n$, $L_n(\widehat h_n)\leq\widehat J_n(\widehat h_n)-C_n+\Delta_n\leq\widehat J_n(h_\eta)-C_n+\Delta_n\leq L_n(h_\eta)+2\Delta_n$.\\
Thus $1\leq L_n(\widehat h_n)/L_n^*\leq1+\eta+2\Delta_n/L_n^*$. As $\Delta_n/L_n^*\xrightarrow{\mathbb{P}}0$, $\eta$ arbitrary, ratio converges to $1$ in probability.\\
\hlt{Structural Commentary}: Uniform risk-equivalence makes cross-validation oracle-optimal.\\
\hlt{Key Insight}: Uniform score error preserves the minimizer's risk.
\end{proof}
\end{theorem}

\begin{method} \label{meth:kde-confidence-band-grid} \hlt{Grid Confidence Band for KDE}
\begin{enumerate}[label=(\roman*)]
\setlength{\itemsep}{0pt}
\item Choose a grid $\mathcal{V}=\{v_1,\ldots,v_N\}$ and set $Y_i(v)\defeq h^{-1}K((v-X_i)/h)$, so $\widehat f_n(v)=\overline Y_n(v)$.
\item Estimate $\operatorname{se}(v)$ by $s(v)/\sqrt n$, where $s^2(v)\defeq(n-1)^{-1}\sum_{i=1}^{n}(Y_i(v)-\overline Y_n(v))^2$.
\item Define $\operatorname{ESS}(v)\defeq\sum_{i=1}^{n}K((v-X_i)/h)/K(0)$, choose approximate $m$ independent grid blocks.
\item Set $q\defeq\Phi^{-1}((1+(1-\alpha)^{1/m})/2)$.
\item Define $\ell(v)\defeq\max\{\widehat f_n(v)-q\,\operatorname{se}(v),0\}$ and $u(v)\defeq\widehat f_n(v)+q\,\operatorname{se}(v)$.
\item The grid-based target is $\mathbb{P}\{\ell(v)\leq f(v)\leq u(v)\ \forall v\in\mathcal{V}\}\approx1-\alpha$.
\end{enumerate}
\end{method}

\subsubsection{Nonparametric Regression}

\begin{definition} \label{def:nadaraya-watson-kernel-regression} \hlt{Nadaraya--Watson Kernel Regression}\\
For data $(x_i,Y_i)$ with $Y_i=r(x_i)+\varepsilon_i$ and $\mathbb{E}\varepsilon_i=0$, the \hlt{Nadaraya--Watson estimator} is $\widehat r(x)\defeq\sum_{i=1}^{n}w_i(x)Y_i$, where $w_i(x)\defeq K((x-x_i)/h)/\sum_{j=1}^{n}K((x-x_j)/h)$.\\
Thus $\widehat r(x)$ is a kernel-weighted local average of nearby responses.
\end{definition}

\begin{proposition} \label{prop:binned-nonparametric-regression-risk} \hlt{Binned Regression Risk}\\
Let $B_j=[(j-1)h,jh)$ for $j<m$ and $B_m=[(m-1)h,1]$, where $h=1/m$.\\
For $x\in B_j$, set $\widehat r_n(x)\defeq\overline Y_j\defeq N_j^{-1}\sum_{i:x_i\in B_j}Y_i$.\\
Assume $x_i$ are approximately regular, $N_j\approx nh$, $\operatorname{Var}(\varepsilon_i)=\sigma^2$, and $\int_0^1(r'(u))^2\,du<\infty$.\\
Then $R(r,\widehat r_n)\approx h^2\int_0^1(r'(u))^2\,du/12+\sigma^2/(nh)$. Hence $h^*\asymp n^{-1/3}$ and $R(r,\widehat r_{h^*})\asymp n^{-2/3}$.
\begin{proof} For $x\in B_j$, $\mathbb{E}\widehat r_n(x)$ is the empirical bin average of $r(x_i)$, approximately the bin average of $r$.\\
Taylor expansion around the bin midpoint gives integrated squared bias $\approx h^2\int_0^1(r'(u))^2\,du/12$.\\
Since $\overline Y_j$ averages $N_j$ independent errors, $\operatorname{Var}(\widehat r_n(x))=\sigma^2/N_j\approx\sigma^2/(nh)$ on $B_j$.\\
Integrating over $x$ preserves variance order, risk proxy is $Q(h)=Ah^2/12+\sigma^2/(nh)$ with $A\defeq\int_0^1(r'(u))^2\,du$.\\
Minimising $Q$ gives $h^*\asymp n^{-1/3}$ and $Q(h^*)\asymp n^{-2/3}$.\\
\hlt{Structural Commentary}: Binned regression is histogram smoothing with response noise.\\
\hlt{Key Insight}: Larger bins reduce noise variance but increase bin-averaging bias.
\end{proof}
\end{proposition}

\begin{theorem} \label{thm:nadaraya-watson-risk-rate} \hlt{Nadaraya--Watson Risk Rate}\\
Assume $(S)$: $r\in C^2$, the design density is positive and smooth, $\operatorname{Var}(\varepsilon\mid X)=\sigma^2<\infty$, and $K$ is second order.\\
If $R(r,\widehat r_h)\approx Q_n(h)\defeq Ah^4+B/(nh)$ with $A,B>0$, then $h^*\asymp n^{-1/5}$ and $R(r,\widehat r_{h^*})\asymp n^{-4/5}$.
\begin{proof} Let $b_h(x)\defeq\mathbb{E}\widehat r_h(x)-r(x)$ and $v_h(x)\defeq\operatorname{Var}(\widehat r_h(x))$. Under $(S)$, Taylor expansion and the second-order moment conditions give $b_h(x)=h^2a(x)+o(h^2)$, so $B_h\defeq\int b_h(x)^2\,dx=Ah^4+o(h^4)$.\\
The kernel window contains order $nh$ observations, hence $V_h\defeq\int v_h(x)\,dx=B/(nh)+o((nh)^{-1})$.\\
Therefore $R(r,\widehat r_h)=B_h+V_h=Q_n(h)+o(h^4+(nh)^{-1})$.\\
Since $Q_n'(h)=4Ah^3-B/(nh^2)$, the optimizer satisfies $h^5=B/(4An)$, hence $h^*\asymp n^{-1/5}$ and $Q_n(h^*)\asymp n^{-4/5}$.\\
\hlt{Structural Commentary}: Risk is smoothing bias plus local-average variance.\\
\hlt{Key Insight}: Balance $h^4$ with $(nh)^{-1}$.
\end{proof}
\end{theorem}

\begin{method} \label{meth:nonparametric-regression-cross-validation} \hlt{Regression Bandwidth Cross-Validation}
\begin{enumerate}[label=(\roman*)]
\setlength{\itemsep}{0pt}
\item Choose $h$ by minimizing leave-one-out score $\widehat J(h)\defeq\sum_{i=1}^{n}(Y_i-\widehat r_{-i}(x_i))^2$, where $\widehat r_{-i}$ omits observation $i$.
\item For a kernel smoother, use $\widehat J(h)\approx\sum_{i=1}^{n}(Y_i-\widehat r(x_i))^2(1-K(0)/\sum_{j=1}^{n}K((x_i-x_j)/h))^{-2}$.
\end{enumerate}
\end{method}

\begin{proposition} \label{prop:adjacent-difference-error-variance-estimator} \hlt{Adjacent-Difference Estimate of Noise Variance}\\
Let $Y_i=r(x_i)+\varepsilon_i$, where $x_1<\cdots<x_n$, $\max_i|x_{i+1}-x_i|\to0$, $r$ is Lipschitz, and $\varepsilon_i$ are independent with $\mathbb{E}\varepsilon_i=0$ and $\operatorname{Var}(\varepsilon_i)=\sigma^2$. Define $\widehat\sigma^2\defeq(2(n-1))^{-1}\sum_{i=1}^{n-1}(Y_{i+1}-Y_i)^2$.\\
Then $\widehat\sigma^2\xrightarrow{\mathbb{P}}\sigma^2$ under the usual averaging conditions for adjacent differences.
\begin{proof} Write $\Delta_iY=Y_{i+1}-Y_i=\Delta_ir+\Delta_i\varepsilon$, with $\Delta_ir=r(x_{i+1})-r(x_i)$ and $\Delta_i\varepsilon=\varepsilon_{i+1}-\varepsilon_i$.\\
Lipschitz continuity and $\max_i|x_{i+1}-x_i|\to0$ give $\max_i|\Delta_ir|\to0$.\\
Also $\mathbb{E}(\Delta_i\varepsilon)^2=2\sigma^2$, so $(2(n-1))^{-1}\sum_i(\Delta_i\varepsilon)^2\xrightarrow{\mathbb{P}}\sigma^2$ by averaging.\\
The cross term and $\sum_i(\Delta_ir)^2$ are $o_{\mathbb{P}}(1)$, hence $\widehat\sigma^2\xrightarrow{\mathbb{P}}\sigma^2$.\\
\hlt{Structural Commentary}: Adjacent differencing cancels a smooth regression signal.\\
\hlt{Key Insight}: Nearby responses differ mainly through independent noise.
\end{proof}
\end{proposition}

\begin{method} \label{meth:kernel-regression-confidence-band} \hlt{Kernel Regression Confidence Band}
\begin{enumerate}[label=(\roman*)]
\setlength{\itemsep}{0pt}
\item Choose a grid $\mathcal{V}$ and define $w_i(v)\defeq K((v-x_i)/h)/\sum_{j=1}^{n}K((v-x_j)/h)$.
\item Estimate noise by $\widehat\sigma^2\defeq(2(n-1))^{-1}\sum_{i=1}^{n-1}(Y_{i+1}-Y_i)^2$.
\item Set $\operatorname{se}(v)\defeq\widehat\sigma(\sum_{i=1}^{n}w_i(v)^2)^{1/2}$.
\item With a grid critical value $q$, define $\ell(v)\defeq\widehat r(v)-q\,\operatorname{se}(v)$ and $u(v)\defeq\widehat r(v)+q\,\operatorname{se}(v)$.
\item The band targets the smoothed regression curve on $\mathcal{V}$.
\end{enumerate}
\end{method}

\begin{remark} \label{rem:curve-confidence-bands-bias-warning} \hlt{Confidence Bands and Bias}\\
Normal bands require negligible bias: if $\widehat\theta$ has bias $b_n$ and standard error $s_n$, nominal coverage needs $b_n/s_n\to0$.\\
For KDE, with $c_1\defeq\int x^2K(x)\,dx$ and $c_2\defeq\int K(x)^2\,dx$, $\widehat f(x)\approx N(f(x)+c_1h^2f''(x)/2,c_2f(x)/(nh))$.\\
At $h\asymp n^{-1/5}$, bias and standard error are both order $n^{-2/5}$, so bands without bias correction target a smoothed density and may undercover $f$.
\end{remark}

\begin{definition} \label{def:additive-nonparametric-regression} \hlt{Additive Nonparametric Regression}\\
For $X=(X_1,\ldots,X_p)$, an \hlt{additive regression model} has $Y=\sum_{j=1}^{p}r_j(X_j)+\varepsilon$.\\
Optional interactions give $Y=\sum_{j=1}^{p}r_j(X_j)+\sum_{j<k}r_{jk}(X_j,X_k)+\varepsilon$.\\
The main-effect model replaces one $p$-dimensional function by $p$ one-dimensional functions.
\end{definition}

\begin{method} \label{meth:additive-regression-backfitting} \hlt{Backfitting}
\begin{enumerate}[label=(\roman*)]
\setlength{\itemsep}{0pt}
\item Initialise component estimates $\widehat r_1,\ldots,\widehat r_p$.
\item For $j=1,\ldots,p$, compute partial residuals $e_i^{(j)}\defeq Y_i-\sum_{s\neq j}\widehat r_s(X_{is})$.
\item Update $\widehat r_j$ by smoothing $e_i^{(j)}$ against $X_{ij}$.
\item Iterate over $j$ until $\widehat r_1,\ldots,\widehat r_p$ converge.
\item The restriction is structural: the target must be well approximated by $\sum_{j=1}^{p}r_j(X_j)$.
\end{enumerate}
\end{method}
